%! Inverse Functions — Unit 1, Lesson 1.6 — Warm-Up
\documentclass[10pt]{article}
\usepackage{precalculus-article}
\usepackage{precalculus-boxes}
\newcommand{\TallMath}[1]{$\displaystyle #1\rule[-1.4em]{0pt}{3.2em}$}

\begin{document}

\pageheader{Unit 1, Lesson 1.6}{Warm-Up}

\vspace{0.12in}

\begin{enumerate}[itemsep=12pt]

\item From Algebra 1: each equation gives you the \textit{answer} and asks for
      the \textit{question}. Solve for $x$.

\begin{work}
  3x + 2 &= 14 \\
      3x &= 12 \\
       x &= 4 \\
  4x - 1 &= 11 \\
      4x &= 12 \\
       x &= 3
\end{work}

\begin{enumerate}[label=(\alph*), itemsep=5pt]
  \item $3x + 2 = 14$ \qquad $x = $ \blank{2cm}
  \item $4x - 1 = 11$ \qquad $x = $ \blank{2cm}
\end{enumerate}

\item From Lesson 1.5: let $f(x) = 3x + 2$ and $g(x) = \dfrac{x - 2}{3}$.
      Name the handoff first, every time.

\begin{work}
  f(g(11)) &= f(3) \\
           &= 11 \\
   g(f(3)) &= g(11) \\
           &= 3
\end{work}

\begin{enumerate}[label=(\alph*), itemsep=5pt]
  \item $f(g(11))$: \; handoff $= $ \blank{1.4cm}, \;
        so $f(g(11)) = $ \blank{1.6cm}
  \item $g(f(3))$: \; handoff $= $ \blank{1.4cm}, \;
        so $g(f(3)) = $ \blank{1.6cm}
  \item Both times you landed back on the number you started with. Do not
        explain it yet --- just write down that it happened.
\end{enumerate}

\item From Lesson 1.1: the table gives every value of a function $q$.

\smallskip
\renewcommand{\arraystretch}{1.4}
\begin{tabular}{|c|c|c|c|c|c|}
\hline
$x$ & 0 & 1 & 2 & 3 & 4 \\
\hline
$q(x)$ & 7 & 3 & 5 & 3 & 1 \\
\hline
\end{tabular}
\smallskip

\begin{enumerate}[label=(\alph*), itemsep=5pt]
  \item $q(2) = $ \blank{2cm}
  \item Now read it backwards: for which input is $q(x) = 7$?
        \quad $x = $ \blank{2cm}
  \item Read $q(x) = 3$ backwards and something goes wrong. The output 3 came
        from \textit{two} inputs: $x = $ \blank{1.4cm} and $x = $
        \blank{1.4cm}.
\end{enumerate}

\end{enumerate}

\end{document}
